\chapter{年度推移と比較の分母}
平成23・24年度は基礎25問、それ以後は基礎30問である。適性15問、応用理学35問は各試験回で同じ登録範囲を使う。出題数の変化を概念の流行と判断する前に分母を確認する。
\begin{longtable}{lrrrr}
試験回 & 基礎 & 適性 & 応用理学 & 合計\\\hline
H23 & 25 & 15 & 35 & 75\\
H24 & 25 & 15 & 35 & 75\\
H25 & 30 & 15 & 35 & 80\\
H26 & 30 & 15 & 35 & 80\\
H27 & 30 & 15 & 35 & 80\\
H28 & 30 & 15 & 35 & 80\\
H29 & 30 & 15 & 35 & 80\\
H30 & 30 & 15 & 35 & 80\\
R01 & 30 & 15 & 35 & 80\\
R01\_retest & 30 & 15 & 35 & 80\\
R02 & 30 & 15 & 35 & 80\\
R03 & 30 & 15 & 35 & 80\\
R04 & 30 & 15 & 35 & 80\\
R05 & 30 & 15 & 35 & 80\\
R06 & 30 & 15 & 35 & 80\\
R07 & 30 & 15 & 35 & 80\\
\end{longtable}
\section{令和3〜7年度の五回と全期間の比較}
後期五回は基礎150、適性75、応用理学175問を分母にする。全期間の構成比と後期の構成比は、同じ主分類で比較する。数ポイントの差だけで単調な増減や次年度の出題を断定しない。
\begin{longtable}{lp{.40\linewidth}rrr}
科目 & 中分類 & 全期間\% & 後期問数 & 後期\%\\\hline
基礎 & エネルギー & 6.2 & 9 & 6.0\\
基礎 & 生物・バイオ & 6.0 & 10 & 6.7\\
基礎 & 科学技術史・社会 & 5.5 & 9 & 6.0\\
基礎 & 材料・金属 & 5.5 & 9 & 6.0\\
基礎 & 環境・持続可能性 & 5.5 & 10 & 6.7\\
基礎 & 微積分・数値解析 & 5.5 & 8 & 5.3\\
適性 & 安全・リスク & 24.6 & 16 & 21.3\\
適性 & コンプライアンス・説明責任 & 22.1 & 17 & 22.7\\
適性 & 技術者倫理・意思決定 & 17.1 & 11 & 14.7\\
適性 & 技術士法・義務責務 & 12.1 & 8 & 10.7\\
適性 & 研究・情報倫理 & 12.1 & 8 & 10.7\\
適性 & 環境・持続可能性 & 7.9 & 9 & 12.0\\
応用理学 & 物理探査 & 12.1 & 21 & 12.0\\
応用理学 & 物理化学 & 8.6 & 16 & 9.1\\
応用理学 & 電磁気・物性 & 7.1 & 14 & 8.0\\
応用理学 & 有機・高分子化学 & 5.9 & 8 & 4.6\\
応用理学 & 応用地質・地盤 & 4.8 & 4 & 2.3\\
応用理学 & 力学 & 4.6 & 9 & 5.1\\
\end{longtable}
再試験を通常年の傾向へ二重に算入する効果、分類境界の変更、主題が同じでも問い方が変わる効果を区別する。ここでは長期構成と後期構成の記述的比較までを行い、有意差や予測確率は計算していない。
